\chapter*{Abstract}\label{chap:00-resume}
\addcontentsline{toc}{chapter}{Abstract}

\paragraph{The question and the answer}
In 2018, at ICLR and then at NeurIPS, Joelle Pineau showed why the published results of deep
reinforcement learning (RL) inspired little trust. We pick out sixteen flaws from her two talks and ask,
for each one, whether the papers from 2018 to 2026 corrected it. Joelle Pineau was heard on what accompanies the result, much less on how the result is judged. The
code link rose from 31\% to 74\% of classical RL papers, the statistical test only from 6\% to 12\%.

\paragraph{Two instruments}
We reread the papers with Pineau's reading grid, extended by criteria of our own: fifty classical RL papers from 2026 against fifty papers from 2018, plus twenty papers per year in
2019, 2020, 2022 and 2024 to date the change. We also rerun the experiments she showed, with four of today's programs on four simulated locomotion
tasks (MuJoCo), ten then twenty random seeds per program and per task.

\paragraph{What the survey finds}
Four practices rose in an established way: the code link, the stated compute, the effect of a
setting and reported uncertainty, the last one less solidly than the other three. The statistical test, held-out evaluation, the tuning of
baselines and the benchmark version did not move detectably. The rise is slow, with no jump after the
checklist NeurIPS has required since 2019.

\paragraph{What the survey cannot say}
Three limits frame these verdicts. The readers are agentic AI models: every \enquote{yes} quotes a sentence of the paper,
which we find again in the PDF, and eight new readers, then a reader from another model
family, find the direction of all four rises. The starting point is our own reading of 2018, not
Pineau's count: she says neither which papers she read nor how she decided, and we find 55\% of papers giving all their
settings where she counted 90\%. We therefore compare our reading of 2018 with our
reading of 2026, both made with the same rules. Finally,
with fifty papers per year, only a rise of 18 to 30 points was detectable: \enquote{no detectable
change} does not mean \enquote{no change}.

\paragraph{The experiments}
Chance fools the eye, sorting fools the test. Between two groups of five runs of one program drawn at
random, Welch's test sees a difference in 1.2\% of draws. Between the five best and the five worst, it
sees one in 31\% of draws. Keeping the best three runs out of ten inflates the mean by about a third.
The $t$ of 9 shown in 2018 probably comes from a computation that ignores the spread between random
seeds. The best program changes with the task, the budget and the physics. At the protocol's budgets, the
leader changes from one task to the next. At equal budget, CrossQ leads on all four tasks. Under a
physics modified by 20\%, PPO JAX, the leader on \task{Hopper}, ends last.
Finally, a published recipe ages: the RL Baselines3 Zoo's recipe for PPO today leaves a robot on \task{Hopper}
standing still, while the agent published with it in 2022 runs. Changing a single setting, the discount
factor, is enough to make today's robot run again.

\paragraph{An exploratory strand on language models}
At a small scale (two models of 0.5 and 1~billion parameters, five random seeds), GRPO
(\emph{Group Relative Policy Optimization}) brings no
established gain. The apparent gain of Qwen, the 0.5-billion model, comes from
shorter answers, cut off less often by the length limit, not from more
correct ones. This strand shows a measurement trap, without judging GRPO.

\paragraph{The short version}
The blog post \href{https://deraison.ai/en/blog/listen-to-joelle-pineau}{\texttt{deraison.ai/en/blog/listen-to-joelle-pineau}}
tells the story of this work in a few pages, for readers who lack the time for this report.

\paragraph{Code and data}
Repository \repo{} (BSD 3-Clause licence): protocols and quoted readings (\file{study/}), training runs
(\file{sim/}), statistics (\file{src/pineau\_rl/}). \file{make verify} regenerates and checks every
result. A protocol written in advance costs little next to a false conclusion.
